\documentclass[12pt, a4paper]{article}

\usepackage{amsmath, amssymb, amsfonts}
\usepackage{float}
\usepackage{graphicx}
\usepackage{listings}
\usepackage{xcolor}
\usepackage[linesnumbered, ruled]{algorithm2e}
\usepackage{subcaption}
\usepackage[hidelinks]{hyperref}
% Each algorithm is written twice: the algorithm2e block (what the PDF typesets)
% and a plain-text Verbatim block right after it (what the web shows — pandoc
% renders Verbatim as a copyable <pre> code block). \excludecomment makes pdflatex
% SKIP the Verbatim blocks, so the PDF shows only the algorithm2e version; pandoc
% conversely drops each algorithm2e to an (empty, CSS-hidden) div and keeps the
% Verbatim. Because pdflatex never sees the Verbatim, its Unicode maths glyphs
% (⟨⟩ ← τ ∈ ∞ …) need no inputenc/literate setup.
\usepackage{comment}
\excludecomment{Verbatim}

\definecolor{codegreen}{rgb}{0,0.6,0}
\definecolor{codegray}{rgb}{0.5,0.5,0.5}
\definecolor{codepurple}{rgb}{0.58,0,0.82}
\hypersetup{linktoc=all}

\lstset{
    language=Python,
    basicstyle=\ttfamily\small,
    keywordstyle=\color{blue},
    commentstyle=\color{codegreen},
    stringstyle=\color{codepurple},
    breaklines=true,
    numbers=left,
    numberstyle=\tiny\color{codegray},
    frame=single,
    showstringspaces=false
}

\title{AI Masterclass Test 1}
\author{}
\date{}


\begin{document}

\section{Uniform Cost Search}
\begin{algorithm}[H]
% \caption{Uniform Cost Search (UCS)}
\SetKwFunction{Fucs}{UCS}
\SetKwProg{Fn}{function}{}{end function}

\Fn{\Fucs{$\mathcal{P} = \langle S, s_0, A, \tau, c, G \rangle$}}{
    $g(s) \leftarrow \infty$ for all $s \in S$\;
    $agenda \leftarrow \{s_0\}$ \tcp*[r]{agenda is a set}
    \While{$agenda \neq \emptyset$}{
        $C \leftarrow \arg \min_{s \in agenda} g(s)$\;
        $s' \leftarrow$ any element of $C$\;
        \If{$s' \in G$}{
            \tcp*[h]{goal found?} \;
            \Return "done"\;
        }
        \For{$a \in A$}{
            \tcp*[h]{expand s} \;
            $s'' \leftarrow \tau(s', a)$\;
            update $g(s'')$ -- replace existing value if smaller\;
            $agenda \leftarrow agenda \cup \{s''\}$\;
        }
    }
}
\end{algorithm}
\begin{Verbatim}
function UCS(P = ⟨S, s₀, A, τ, c, G⟩):
    g(s) ← ∞  for all s ∈ S
    agenda ← {s₀}                     # agenda is a set
    while agenda ≠ ∅:
        C ← arg min_{s ∈ agenda} g(s)
        s' ← any element of C
        if s' ∈ G:                    # goal found?
            return "done"
        for a ∈ A:                    # expand s
            s'' ← τ(s', a)
            update g(s'')  (replace existing value if smaller)
            agenda ← agenda ∪ {s''}
\end{Verbatim}

\begin{itemize}
    \item We have the cost function
    \[c: S \times A \times S \rightarrow \mathbb{R}\]
    \item UCS involves keeping cost of the path to any node. Let this be $g(s) \in \mathbb{R}$.
    \item UCS utilises a priority queue, where the priority is the path cost.
    \item UCS expands cheapest nodes first (in order of incerasing path cost).
    \item Always finds optimal path if costs are non-negative ($\implies$ path costs either stay the same or increase). 
    If UCS found the goal with cost $C$,
    there can't be another path with cost $C' < C$, since UCS would already have expanded the path,
    by definition of the priority queue.
    \item UCS becomes BFS when all step costs are equal.
\end{itemize}

\begin{figure}[H]
    \centering
    % --- Left Column (Image) ---
    \begin{minipage}{0.45\textwidth}
        \centering
        \includegraphics[width=\textwidth]{images/ucs-eg.png}
        \caption{Handwritten diagram}
    \end{minipage}
    \hfill % Adds horizontal space between the two columns
    % --- Right Column (Text/Code) ---
    \begin{minipage}{0.45\textwidth}
    Execution Example
    \begin{itemize}
    \item Successors are Rimnicu (80) and Fagaras (99)
    \item Rimnicu is selected for expansion
    \item Successor is Pitesi ($80 + 97 = 177$)
    \item Least cost is now Fagaras
    \item Successor is Bucharest ($99 + 211 = 310$)
    \item Pitesi now selected for expansion
    \item Bucharest updated ($80 + 97 + 101 = 278$)
    \item Bucharest selected for expansion; goal reached with optimal path
    \end{itemize}    
    \end{minipage}
\end{figure}

\end{document}
